\begin{python0}
from solutions import *; clear()
\end{python0}
\sectionthree{Friends}

You can let \textit{operator<<} have access to the
private members of the object by making your
\textit{operator<<} a \textbf{friend} in your class. We
will re-visit friends later.

\begin{consolethree}[escapeinside=||]
#include <iostream>

#ifndef INT_H
#define INT_H

class C
{
public:
    Int(int x0 = 0)
        : x_(x0)
    {}
    int get() const { return x_; }
    |{\bf friend}| std::ostream & operator<<(std::ostream &, const Int &);
private:
    int x_;
};

#endif
\end{consolethree}

\begin{consolethree}[escapeinside=||]
#include "Int.h"

std::ostream & operator<<(std::ostream & cout, const Int & c)
{
    cout << |{\bf c.x\_}|;
    return cout;
}
\end{consolethree}

\begin{consolethree}[escapeinside=||]
#include "Int.h"

int main()
{
    Int c(42),d(43);
    std::cout << c << "," << d << "\n";
    return 0;
}
\end{consolethree}

In general, if the header of a function \textit{f} is a friend in a class
\textit{C}, then \textit{f} has access to things in private section of class
\textit{C}:

\begin{consolethree}[escapeinside=||]
class C
{
public:
    |{\bf friend void f(C \& c);}|
private:
    int x_;
    void m() { x++; }
}

void f(C & c)
{
    c.x_ = 0; // f can access c.x_ (private)
    c.m(); // f can access c.m (private)
}
\end{consolethree}

Do NOT declare too many friends in a class -- this breaks the principle
of information hiding and makes the class harder to maintain in case you
need private members variables or methods in the future.

Of course friends can be completely avoided.


